% ============================================================
% Section 12: Mid-Late 20th Century (Slides 84-91)
% ============================================================
\section{Mid-Late 20th Century}

% ----------------------------------------------------------
% Slide 84: Section Divider
% ----------------------------------------------------------
{
\setbeamercolor{background canvas}{bg=digitalBlue}
\begin{frame}[plain]
  \centering
  \vspace{2cm}
  \textcolor{white}{\Huge\bfseries Mathematics Goes Digital}\\[1em]
  \textcolor{white!80}{\large 1950--2000 CE}\\[1.5em]
  \textcolor{white!60}{\normalsize Shannon \(\cdot\) Lorenz \(\cdot\) Mandelbrot \(\cdot\) Wiles}\\[0.3em]
  \textcolor{white!60}{\normalsize Cryptography \(\cdot\) Computer-Assisted Proof}

  \note{
    This half-century saw mathematics go from a purely abstract discipline to
    the hidden engine of the digital world.
    Information theory, chaos, fractals, cryptography, and computer-assisted proof
    all emerged in this period.
  }
\end{frame}
}

% ----------------------------------------------------------
% Slide 85: Claude Shannon and Information Theory
% ----------------------------------------------------------
\begin{frame}{Claude Shannon --- The Bit and Information Theory}

  \begin{columns}[T]
    \begin{column}{0.48\textwidth}
      \begin{personbox}[Claude Elwood Shannon (1916--2001)]{digitalBlue}
        \textbf{Origin:} Petoskey, Michigan, USA; worked at Bell Labs and MIT\\[3pt]
        \textbf{Key work:} Founded \alert{information theory}
        (``A Mathematical Theory of Communication,'' 1948). Defined the
        \textbf{bit} as the fundamental unit of information. Showed that
        all communication can be reduced to binary digits.
      \end{personbox}

      \vspace{0.4em}
      \begin{insightbox}
        Every text message, every video stream, every bit of data on the
        internet exists because of Shannon's mathematics.
      \end{insightbox}
    \end{column}

    \begin{column}{0.49\textwidth}
      % TikZ: Shannon communication model
      \visualplaceholder[5cm]{Blocks}
    \end{column}
  \end{columns}

  \note{
    ``Every text message, every video stream, every bit of data on the internet
    exists because of Shannon's mathematics.''
    Shannon also proved fundamental limits: the channel capacity theorem tells us
    the maximum rate at which information can be reliably transmitted.
    Fun fact: Shannon was also a juggler and built a flame-throwing trumpet.
  }
\end{frame}

% ----------------------------------------------------------
% Slide 86: Edward Lorenz and Chaos Theory
% ----------------------------------------------------------
\begin{frame}{Edward Lorenz --- The Butterfly Effect}

  \begin{columns}[T]
    \begin{column}{0.45\textwidth}
      \begin{personbox}[Edward Norton Lorenz (1917--2008)]{digitalBlue}
        \textbf{Origin:} West Hartford, Connecticut, USA; worked at MIT\\[3pt]
        \textbf{Key work:} Discovered \alert{deterministic chaos} (1963)
        while running a weather simulation. Tiny rounding differences in
        initial conditions produced wildly different outcomes ---
        the \textbf{butterfly effect}. Even perfectly deterministic systems
        can be practically unpredictable.
      \end{personbox}

      \vspace{0.3em}
      \begin{insightbox}
        ``Does the flap of a butterfly's wings in Brazil set off a tornado
        in Texas?'' Not magic --- mathematics.
      \end{insightbox}
    \end{column}

    \begin{column}{0.52\textwidth}
      % TikZ: Lorenz attractor approximation
      \visualplaceholder[5cm]{Approximate the butterfly shape with parametric curves}
    \end{column}
  \end{columns}

  \note{
    ``He discovered that the flap of a butterfly's wings in Brazil could, in
    theory, set off a tornado in Texas. Not magic --- mathematics.''
    Lorenz found this by accident: he re-entered initial conditions rounded to
    3 decimal places instead of 6. The weather simulation diverged completely.
    This showed that long-term weather prediction has fundamental limits ---
    even with perfect equations.
    The Lorenz attractor is one of the iconic images of mathematics.
  }
\end{frame}

% ----------------------------------------------------------
% Slide 87: Four-Color Theorem -- Computers Enter Mathematics
% ----------------------------------------------------------
\begin{frame}{The Four-Color Theorem (1976) --- Computers Enter Mathematics}

  \begin{columns}[T]
    \begin{column}{0.5\textwidth}
      \begin{personbox}[Kenneth Appel (1932--2013) \& Wolfgang Haken (1928--2022)]{digitalBlue}
        \textbf{Proved at:} University of Illinois, 1976\\[3pt]
        Any map can be coloured with at most \alert{four colours} such that
        no two adjacent regions share a colour. Appel and Haken used a
        computer to check reducible configurations --- the \textbf{first
        major theorem proved with essential computer assistance}.
      \end{personbox}

      \vspace{0.3em}
      \begin{insightbox}
        Is a proof that no human can fully verify still a ``real'' proof?
        The debate continues today.
      \end{insightbox}

      \vspace{0.2em}
      {\scriptsize\itshape The UIUC math department used a postmark reading
      ``Four Colors Suffice'' to celebrate.}
    \end{column}

    \begin{column}{0.47\textwidth}
      % TikZ: Map coloring example
      \visualplaceholder[5cm]{A simple map requiring 4 colors}
    \end{column}
  \end{columns}

  \note{
    ``For the first time, a computer was essential to proving a mathematical theorem.
    Some mathematicians refused to accept it. The debate about computer-assisted
    proofs continues today.''
    \verifytag{Reducible configurations count: sometimes cited as 1,476 or 1,936.
    Verify exact number. UIUC postmark detail to be verified.}
    The proof was later simplified and verified by other computer programs
    (Robertson et al., 1997).
  }
\end{frame}

% ----------------------------------------------------------
% Slide 88: Benoit Mandelbrot and Fractals
% ----------------------------------------------------------
\begin{frame}{Benoit Mandelbrot --- Fractals and Hidden Patterns}

  \begin{columns}[T]
    \begin{column}{0.42\textwidth}
      \begin{personbox}[Benoit B.\ Mandelbrot (1924--2010)]{digitalBlue}
        \textbf{Origin:} Warsaw, Poland (Lithuanian-Jewish origin);
        worked in France and at IBM/Yale in the USA\\[3pt]
        \textbf{Key work:} Founded \alert{fractal geometry}. The Mandelbrot
        set (1980) revealed infinite complexity from simple equations.
        Natural phenomena --- coastlines, mountains, blood vessels,
        financial markets --- have fractal structure.
      \end{personbox}

      \vspace{0.3em}
      \begin{insightbox}
        The messiest, most irregular things in nature follow hidden
        mathematical patterns.
      \end{insightbox}
    \end{column}

    \begin{column}{0.55\textwidth}
      % TikZ/pgfplots: Mandelbrot set approximation
      \visualplaceholder[5cm]{We use a simplified visual representation since true Mandelb}
    \end{column}
  \end{columns}

  \note{
    ``He showed that the messiest, most irregular things in nature --- coastlines,
    clouds, mountain ranges --- follow hidden mathematical patterns.''
    Mandelbrot coined the word ``fractal'' from Latin \emph{fractus} (broken).
    The Mandelbrot set contains infinite detail at every zoom level, yet is
    defined by one of the simplest equations in mathematics.
    Fun fact: ``The B.\ in Benoit B.\ Mandelbrot stands for Benoit B.\ Mandelbrot''
    (a fractal joke attributed to him).
  }
\end{frame}

% ----------------------------------------------------------
% Slide 89: Cryptography -- The Mathematics of Secrecy
% ----------------------------------------------------------
\begin{frame}{Cryptography --- The Mathematics of Secrecy}

  \begin{columns}[T]
    \begin{column}{0.48\textwidth}
      {\small Public-key cryptography, invented in the 1970s, made secure
      internet communication possible. Based on number theory ---
      specifically the difficulty of \textbf{factoring large numbers}.}

      \vspace{0.4em}
      \begin{personbox}[Diffie--Hellman Key Exchange (1976)]{digitalBlue}
        \textbf{Whitfield Diffie} (b.\,1944, USA)\\
        \textbf{Martin Hellman} (b.\,1945, USA)\\[3pt]
        Co-invented public-key cryptography.
      \end{personbox}

      \vspace{0.3em}
      \begin{personbox}[RSA Algorithm (1977)]{digitalBlue}
        \textbf{Ron Rivest} (b.\,1947, USA)\\
        \textbf{Adi Shamir} (b.\,1952, Israel)\\
        \textbf{Leonard Adleman} (b.\,1945, USA)\\[3pt]
        Multiply two large primes --- easy.\\
        Factor the product --- nearly impossible.
      \end{personbox}
    \end{column}

    \begin{column}{0.49\textwidth}
      % TikZ: RSA concept diagram
      \visualplaceholder[5cm]{Alice and Bob}
    \end{column}
  \end{columns}

  \note{
    ``Every time you see the padlock icon in your browser, every time you send a
    message on Signal or WhatsApp --- you are relying on number theory from the
    1970s. The math that Euler and Fermat did for fun now protects your privacy.''
    This is a perfect example of ``useless'' pure mathematics becoming essential
    applied mathematics centuries later.

    NAME THE MOTIF: ``We keep seeing this pattern: mathematics done for pure
    curiosity -- Euler playing with prime numbers in the 1700s, Apollonius
    slicing cones in 200 BCE -- turning out to be essential centuries later.
    The math you do in class that seems `useless'? Give it time. Euler never
    imagined online banking. Apollonius never imagined satellite orbits. The
    math YOUR generation finds beautiful today will solve problems that do
    not exist yet.''
  }
\end{frame}

% ----------------------------------------------------------
% Interactive Moment: Encryption Poll
% ----------------------------------------------------------
\begin{frame}{Quick Poll}

  \begin{center}
    \vspace{1cm}
    {\Huge\bfseries\textcolor{digitalBlue}{Raise your hand if you have}}\\[0.6em]
    {\Huge\bfseries\textcolor{digitalBlue}{used encryption today.}}\\[2em]

    \pause

    {\Large\textcolor{accentOrange}{\bfseries Answer: \underline{everyone}.}}\\[1em]

    \visualplaceholder[5cm]{Diagram}
  \end{center}

  \note{
    Interactive moment --- get hands up first, then reveal the answer.
    Point out: ``The math behind that encryption is number theory from the 1970s.
    Ancient number theory, done for pure intellectual curiosity by Euler and Fermat,
    now protects billions of people every day.''
    Budget: 1--2 minutes for this interaction.

    MIRROR MOMENT: After the reveal (``everyone used encryption today''),
    deliver slowly: ``Every time you unlock your phone, you are using number
    theory from the 1700s. Every time you scroll social media, you are using
    linear algebra. The mathematics we have been studying today is not in a
    museum -- it is in your pocket. And the mathematics YOUR generation will
    create? It is not written yet.'' [Let this land. This is where the
    audience realises they are part of the relay.]
  }
\end{frame}

% ----------------------------------------------------------
% Slide 90: Andrew Wiles and Fermat's Last Theorem
% ----------------------------------------------------------
\begin{frame}{Andrew Wiles --- Fermat's Last Theorem Proved}

  \begin{columns}[T]
    \begin{column}{0.5\textwidth}
      \begin{personbox}[Sir Andrew John Wiles (b.\,1953)]{digitalBlue}
        \textbf{Origin:} Cambridge, England; worked at Princeton\\[3pt]
        \textbf{Key work:} Proved \alert{Fermat's Last Theorem} in 1995
        (after finding and fixing a gap in his 1993 proof). Used deep
        connections between elliptic curves and modular forms
        (the Taniyama--Shimura--Weil conjecture, now the modularity theorem).
        Closed a \textbf{358-year-old problem}.
      \end{personbox}

      \vspace{0.3em}
      \begin{insightbox}
        He worked in secret for seven years. When he announced the proof,
        he famously choked up with emotion.
      \end{insightbox}
    \end{column}

    \begin{column}{0.47\textwidth}
      % Fermat's Last Theorem
      \visualplaceholder[5cm]{The simple statement}
    \end{column}
  \end{columns}

  \note{
    ``He worked in secret for seven years. When he announced the proof, he famously
    choked up with emotion at the end of his lecture.''
    The proof connects three distant areas of mathematics: number theory, algebraic
    geometry, and modular forms. It required proving a special case of the
    Taniyama--Shimura--Weil conjecture.
    Wiles was awarded a special Silver Plaque by the IMU (he was over the Fields
    Medal age limit) and later the Abel Prize (2016).
  }
\end{frame}

% ----------------------------------------------------------
% Slide 91: 20th Century Overview
% ----------------------------------------------------------
\begin{frame}{The 20th Century --- An Overview}

  \begin{center}
    \visualplaceholder[5cm]{Bottom summary}
  \end{center}

  \note{
    Quick recap. Point to each card briefly.
    Key message: mathematics went from purely abstract to the hidden infrastructure
    of modern civilisation.
    Transition: ``Turing's theoretical machine became real. And now, machines are
    beginning to DO mathematics.''
  }
\end{frame}
